\documentclass[11pt,a4paper]{article}
\usepackage[margin=2cm,top=1.95cm,bottom=1.95cm]{geometry}
\usepackage[fontset=fandol]{ctex}
\usepackage{setspace}\AtBeginDocument{\singlespacing}
\usepackage{amsmath,amssymb,bm,booktabs,graphicx,pgfplots,caption,hyperref,array,xcolor,float}
\pgfplotsset{compat=1.18}\usepgfplotslibrary{groupplots}
\pgfplotsset{table/search path={figs/}}
\usetikzlibrary{positioning,arrows.meta,calc}
\captionsetup{font=small,labelfont=bf}\hypersetup{hidelinks,pdftitle={从湍流级联到A股分钟收益},pdfauthor={方洪威，戴晴琴}}
\newcolumntype{L}[1]{>{\raggedright\arraybackslash}p{#1}}
\newcommand{\chart}[4]{\begin{figure}[H]\centering\resizebox{#1\linewidth}{!}{\input{figs/#2.tex}}\caption{#3}\label{#4}\end{figure}}
\setlength{\textfloatsep}{10pt}\setlength{\floatsep}{8pt}\setlength{\intextsep}{9pt}
\title{\textbf{从湍流级联到 A 股分钟收益}\\[4pt]\large 分层波动模型的迁移与逐日预测}
\author{方洪威 \qquad 戴晴琴}\date{}
\raggedbottom
\begin{document}
\maketitle
\section{摘要与物理问题定义}
\subsection*{摘要}
同一只股票可以连续几个小时交投平静，也可以在几分钟内大幅波动。把这些时期的收益放在一起，得到的分布在零点附近较尖，两侧也比高斯分布更厚。湍流的分层模型为这种现象提供了一种解释：短时涨跌的幅度受到变化较慢的方差背景调节。本文从流体中的尺度传递出发，引入用于外汇收益的分层随机方程，再用八只沪市股票的一分钟行情构造五分钟收益，检验这一解释可以走多远。历史交易日按时间以 $8:1:1$ 划分，另用 2026 年行情检验跨期表现。

结果分为两部分。首先，同一时段内两个相邻、各长 40 分钟的窗口往往同时活跃。中国平安在历史测试段与 2026 年的窗口波动秩相关分别为 0.379 和 0.450。这个时间特征支持引入波动背景，但收益和局部方差的联合分布对层数的约束较弱。按股票、采样频率和预处理区分的 12 组拟合中只有 4 组在已测层数范围内出现内部最低点。其次，用前几个交易日的实际波动更新背景，可以改善后续五分钟收益的分布预测。与训练样本相同的静态八层模型相比，条件八层在中国平安两个评价时期的平均对数评分分别提高 0.2148 和 0.0432 纳特/条；八只股票的比较方向一致。主例中，改用每天更新的背景所带来的差异，大于在相同更新规则下由单层增加至八层的差异。

\subsection{研究问题}
价格曲线本身看不出有几层背景。我们先问，短段波动是否具有可重复观察的持续性；再问，静态分布能把层数分辨到什么程度；最后检验，在开盘前已经知道的条件下，哪些信息有助于估计当天五分钟收益的分布。
\chart{.96}{study_map}{本文的论证顺序。第一步是物理类比，中间两步给出可计算、可观测的量，最后一步由后续行情评价。}{fig:map}

\clearpage
\section{模型与算法架构}
\subsection{湍流中的能量如何跨越尺度}
考虑充分发展的三维湍流。外力在大尺度上输入动能，惯性作用把能量逐步输送给较小尺度的运动，最终由黏性耗散为热。一个涡在流体中移动的速度包含大尺度整体运动；要分辨尺度 $\ell$ 上的起伏，通常考察相距 $\ell$ 的两点的纵向速度差
\[
 \delta_\ell u=[\mathbf u(\mathbf r+\ell\mathbf e)-\mathbf u(\mathbf r)]\cdot\mathbf e.
\]
其中 $\mathbf e$ 为取样方向。若 $\ell$ 处在注入与耗散尺度之间，且惯性区主要由平均单位质量能量通量 $\varepsilon$ 控制，量纲估计给出典型速度差 $(\varepsilon\ell)^{1/3}$。这里 $\varepsilon$ 的量纲为 $\mathsf{L}^2\mathsf{T}^{-3}$。这一估计给出典型幅度。局部能量传递强弱并不均匀，固定 $\ell$ 下也会出现成串的剧烈起伏，这称为间歇性。分层随机理论关心的正是背景强度的这种起伏\cite{macedo2017,vasconcelos2024}。

\chart{.73}{energy_cascade}{三维湍流能量传递的尺度示意。每个圆表示一段尺度上的运动，箭头表示能量的平均传递方向。}{fig:cascade}

\subsection{从速度差到外汇收益}
Vasconcelos 等\cite{vasconcelos2024}把空间尺度上的能量级联类比为交易时间尺度之间的信息影响。持仓者在较长时间内作出的调整，可能改变短时交易所处的波动环境。设 $S(t)$ 为货币对在 $t$ 时的汇率，观测量是
\[
 x_\tau(t)=\ln S(t+\tau)-\ln S(t).
\]
两点的对数报价之差与速度增量在模型中承担相似的角色。$\ell$ 是空间距离，$\tau$ 是采样间隔；流体可以测量能量通量，金融模型中的 $v_i$ 则是从价格间接推断的方差。收益统计可以检验波动强弱的关联，却不能凭自身确定信息的传播方向。
\begin{table}[H]\centering\small
\caption{物理图像与可测量量的对应。$\mathsf L$、$\mathsf T$ 分别为长度与时间量纲，$1$ 表示无量纲。}
\begin{tabular}{L{3.1cm}L{5.2cm}L{6.2cm}}\toprule
比较对象 & 流体湍流 & 外汇及股票模型\\\midrule
观测起伏 & 速度增量 $\delta_\ell u$，$\mathsf L\mathsf T^{-1}$ & 报价的对数增量 $r_t$，$1$\\
尺度顺序 & 注入尺度、惯性区、耗散尺度 & 慢市场环境、较快方差背景、短时收益\\
局部强弱 & 局部能量传递或耗散强度随时空变化 & 方差 $v_i$ 随时间变化，以滚动方差估计末层背景\\
可检验的现象 & 小尺度增量出现间歇性强起伏 & 强波动时段的相邻窗口是否也更活跃\\
\bottomrule\end{tabular}\end{table}
论文使用 2012 年 11 月至 2023 年 4 月的小时报价，每个货币对约 6.4 万个点。按收益与局部方差的联合拟合误差，欧元兑美元、美元兑日元、英镑兑美元分别选出 5、2、1 层。在流体中，雷诺数越大，注入尺度与耗散尺度通常分得越开。三组货币对的成交量排序也恰与拟合层数一致；作者由此提出市场规模可能与层数有关。成交量与雷诺数之间没有已知的换算关系。这三个数是各自样本的估计结果；迁移到股票时，需要重新估计层数和模型参数。

\clearpage
\subsection{外汇模型的方程与可计算的密度}\label{sec:model}
模型将给定采样间隔的收益记为 $x$，将从慢到快排列的方差背景记为 $v_0,v_1,\ldots,v_N$。最慢背景 $v_0$ 暂取常数。每层既受上一层牵引，又受本层随机扰动：
\begin{align}
 \mathrm dv_i&=-\gamma_i(v_i-v_{i-1})\,\mathrm dt+\kappa_i v_i\,\mathrm dW_i,\qquad i=1,\ldots,N, \label{eq:hier}\\
 \mathrm dx&=-\mu x\,\mathrm dt+\sqrt{2\mu v_N}\,\mathrm dW_0. \label{eq:fast}
\end{align}
各 $W_i$ 为独立的标准布朗运动，方程取伊藤解释。$\gamma_i$ 和 $\mu$ 的量纲均为 $\mathsf T^{-1}$，$\kappa_i$ 为 $\mathsf T^{-1/2}$。当 $v_i>v_{i-1}$ 时，回复项把本层方差拉低；$\kappa_i v_i$ 则使高方差状态承受更大的绝对扰动。近似条件为 $\gamma_1\ll\cdots\ll\gamma_N\ll\mu$：在快速变量调整的时间内，上一层变化很小。本文采用 $v_N=\operatorname{Var}(x\mid v_N)$ 的记号。按标准布朗运动的归一化，文献\cite{vasconcelos2024}式(1)的噪声幅度写作 $\sqrt{\mu v_N}$，会给出 $v_N/2$；式\eqref{eq:fast}使用 $\sqrt{2\mu v_N}$，与该文式(3)及实际用于拟合的条件高斯密度一致。

固定 $v_{i-1}=u$，令定态概率流为零，可得
\[
 \gamma_i(u-v)f(v\mid u)=\tfrac12\kappa_i^2\,\partial_v[v^2 f(v\mid u)].
\]
定义 $\beta_i=2\gamma_i/\kappa_i^2$，积分并归一化得到逆伽马密度
\begin{equation}
 f(v\mid u)=\frac{(\beta_i u)^{\beta_i+1}}{\Gamma(\beta_i+1)}
 v^{-\beta_i-2}\exp(-\beta_i u/v),\qquad v>0,\quad \mathbb E[v\mid u]=u. \label{eq:kernel}
\end{equation}
因而每层改变方差的分散程度，同时保持给定上一层时的均值。设各层共享一个形状参数 $\beta$，令 $Y_i=v_i/v_{i-1}$，则准静态构造中的乘子相互独立，且 $v_N=v_0\prod_iY_i$。逐层积分给出末层密度 $f_N$，再对条件高斯取平均：
\begin{equation}
 f_i(v)=\int_0^\infty f(v\mid u) f_{i-1}(u)\,\mathrm du,\qquad
 P_N(x)=\int_0^\infty\frac{e^{-x^2/(2v)}}{\sqrt{2\pi v}}f_N(v)\,\mathrm dv,\quad
 f_0(u)=\delta(u-v_0). \label{eq:mix}
\end{equation}
同样的 $x$ 在平静和活跃时段可能有很不一样的概率密度。小方差状态把权重集中在零附近，大方差状态使大幅涨跌更常见。单层时可以直接计算
\[
 P_1(x)=\frac{\Gamma(\beta+3/2)}{\sqrt{2\pi\beta v_0}\,\Gamma(\beta+1)}
 \left(1+\frac{x^2}{2\beta v_0}\right)^{-\beta-3/2}.
\]
这是方差为 $v_0$ 的对称 $t$ 分布。单层已经能够产生厚尾，所以判断多层是否有用，必须同时观察局部方差的分布，并与单层作样本外比较。


\paragraph{方程中哪些量能由价格估计？}
在这一近似下，模型的计算顺序可以写为
\[
 \underbrace{v_0}_{\text{历史方差}}\;
 \xrightarrow[\text{慢}]{\,Y_1\,}v_1
 \xrightarrow{\,Y_2\,}\cdots
 \xrightarrow[\text{快}]{\,Y_N\,}v_N
 \xrightarrow{\ x\mid v_N\ }\underbrace{x}_{\text{对数收益}} .
\]
价格给出最右侧的 $x$。中间各层都是隐变量。只观察一个时点的收益分布时，$\gamma_i$ 与 $\kappa_i$ 以比值 $\beta_i=2\gamma_i/\kappa_i^2$ 出现，二者不能分别估计。共享 $\beta$ 让每个待选层数只有一个形状参数可调；这项约束并未假定所有层的回复时间都相同。

\clearpage
\subsection{把外汇报价换成股票分钟行情}\label{sec:data}
在股票市场，报价变量换为同一只股票的成交价格 $C$。两次报价的比值消去价格单位，因此股票也可使用对数收益。两种市场最直接的差别在交易时钟。外汇研究使用一小时报价；沪市一天分上午、下午两段，午休前后的两个相邻记录可能相隔一个半小时。我们在每一段内用一分钟记录组成五分钟区间，将收盘记录 15:00 并入最后一区间。每个区间须具备至少 80\% 的必需分钟标签，只对相邻有效区间的收盘价取对数比。完整日共 48 个收盘价、46 个五分钟收益。五分钟也是本文的主采样间隔；一分钟和十五分钟另作敏感性检验。

\subsubsection{股票范围与选样}
2017--2019 年共 731 个沪市交易日，1,573 只股票合计 996,288 个有记录的股票日。图\ref{fig:universe}按两种缺失情况展示全体股票。先要求前 584 个交易日中至少 95\% 有记录，再要求有记录的日期中至少 90\% 具备完整分钟标签，且日成交额中位数为正。完整日要求 09:30--11:29、13:00--14:56 的 237 个核心分钟及 15:00 收盘记录齐备。缺记录日按零成交额计入中位数。158 只合格股票中按该中位数取前八，中国平安排第一，作为后续图示的主例。
\chart{.98}{universe_selection}{左图每点为一只股票，横轴衡量有多少天留下记录，纵轴衡量记录齐全的日期占比。虚线为入选门槛。右图按训练期成交额中位数排序，橙点为前八名。}{fig:universe}
\begin{table}[H]\centering\small
\caption{入选的八只股票。排序只依据训练期日成交额中位数；完整日比例的分母为有记录日数，覆盖率的分母为584个名义训练日。}
\label{tab:selected_stocks}
\begin{tabular}{rlcrrrr}\toprule
排名 & 股票 & 证券代码 & 有记录日 & 覆盖率(\%) & 完整日(\%) & 成交额(亿元)\\\midrule
1 & 中国平安 & 601318.SH & 584 & 100.00 & 100.00 & 39.17\\
2 & 贵州茅台 & 600519.SH & 584 & 100.00 & 100.00 & 22.55\\
3 & 中信证券 & 600030.SH & 584 & 100.00 & 98.29 & 15.62\\
4 & 伊利股份 & 600887.SH & 582 & 99.66 & 99.66 & 13.46\\
5 & 招商银行 & 600036.SH & 584 & 100.00 & 100.00 & 12.44\\
6 & 华友钴业 & 603799.SH & 583 & 99.83 & 98.11 & 12.35\\
7 & 方大炭素 & 600516.SH & 583 & 99.83 & 99.49 & 11.01\\
8 & 中国建筑 & 601668.SH & 584 & 100.00 & 100.00 & 10.75\\
\bottomrule\end{tabular}\end{table}

这个选择优先保证日内数据完整、成交活跃。后面的八股比较检验同一种模型在这些股票上是否给出相近结论，完整股票清单随数据附录提供。

\clearpage
\subsubsection{从五分钟价格读出模型变量}
令 $t$ 表示同一时段内第 $t$ 个五分钟区间。外汇模型中的对数增量在这里对应 $r_t$，用于拟合的输入再写为训练波动单位：
\begin{equation}
 r_t=\ln(C_t/C_{t-1}),\quad
 s_{\rm tr}^2=\frac{\sum_{t\in{\rm tr}}(r_t-\bar r_{\rm tr})^2}{n_{\rm tr}-1},
 \qquad x_t=r_t/s_{\rm tr}. \label{eq:returns}
\end{equation}
每只股票单独计算 $s_{\rm tr}$，验证、测试和 2026 年均沿用它。计算标准差时减去均值，输入模型的收益保留原来的零点。以中国平安为例，2017 年 1 月 3 日相邻收盘价由 35.43 元降至 35.40 元，$r_t=-0.00084710$。以 $s_{\rm tr}=0.00221129$ 换算，$x_t=-0.3831$，即跌幅约为训练期五分钟波动标准差的 0.38 倍。

方差背景 $v_N$ 无法从单笔价格直接读出。我们在每个交易时段内取最近 $M$ 个连续收益，计算
\begin{equation}
 \widehat v_{t,M}=\frac{1}{M-1}\sum_{j=0}^{M-1}(x_{t-j}-\bar x_{t,M})^2,\qquad
 \bar x_{t,M}=M^{-1}\sum_{j=0}^{M-1}x_{t-j}. \label{eq:local}
\end{equation}
\chart{.83}{model_data_overview}{中国平安训练段最早的五个完整交易日。自上而下是五分钟收盘价、归一化收益和 $M=8$ 的滚动方差。午休、隔夜均断开；每个新时段先积累八个收益，再画第一个方差点。五天共得到 240 个收盘价、230 个收益和 160 个方差估计。}{fig:observations}

\clearpage
\subsection{数据划分与计算顺序}
历史日期按先后以 $8:1:1$ 划分，名义训练、验证、测试段分别为 584、73、74 日。每段末尾留出 5 个交易日作间隔，股票实际缺失日另计。保留的训练段至 2019 年 5 月 22 日，验证段为 5 月 30 日至 9 月 3 日，测试段为 9 月 11 日至 12 月 24 日。筛选样本股票只看名义训练段，估计尺度、分箱、窗口和模型参数只看保留训练段。
\begin{table}[H]\centering\small
\caption{八股有效收益样本与训练期特征。各样本列为“收益条数/交易日数”；标准差以原始对数收益的百分数表示，零收益比例未作平滑。}
\label{tab:eight_counts}
\begin{tabular}{lrrrrrr}\toprule
股票 & 训练 & 验证 & 测试 & 2026年 & 标准差(\%) & 零收益(\%)\\\midrule
中国平安 & 26\,634/579 & 3\,128/68 & 3\,174/69 & 7\,360/160 & 0.2211 & 5.32\\
贵州茅台 & 26\,634/579 & 3\,128/68 & 3\,174/69 & 7\,360/160 & 0.2290 & 1.66\\
中信证券 & 26\,174/569 & 3\,128/68 & 3\,128/68 & 7\,360/160 & 0.2498 & 14.95\\
伊利股份 & 26\,450/575 & 3\,128/68 & 3\,174/69 & 7\,360/160 & 0.2695 & 9.21\\
招商银行 & 26\,634/579 & 3\,128/68 & 3\,174/69 & 7\,360/160 & 0.2423 & 9.28\\
华友钴业 & 26\,128/568 & 3\,128/68 & 3\,174/69 & 2\,070/45 & 0.4511 & 3.62\\
方大炭素 & 26\,542/577 & 3\,128/68 & 3\,174/69 & 7\,360/160 & 0.3953 & 10.91\\
中国建筑 & 26\,634/579 & 3\,128/68 & 3\,174/69 & 7\,360/160 & 0.2293 & 29.73\\
\bottomrule\end{tabular}\end{table}

2026 年 1 月 5 日至 8 月 31 日另作跨年检验。华友钴业从 6 月 30 日才有记录，单列为 45 日的短样本。2020 年下半年资料只覆盖入选八股中的两只，2021--2025 年缺少可用分钟记录，跨年检验因此另起一段。按既定的 60 个完整交易日、5,000 条收益的双重门槛，2019 年五分钟留出段均属于短样本。
\chart{.97}{transfer_architecture}{从同一批行情进入静态模型和条件模型的估计流程。$M$ 控制观测方差的窗口，$N$ 控制模型中的隐含层数。它们回答不同问题。}{fig:architecture}


\clearpage
\subsection{静态模型如何估计窗口、层数和形状}
\subsubsection{先选择观测窗口}
短窗口跟随波动快，方差估计也更易受单次涨跌影响。长窗口可以平滑误差，却可能把两种背景混在一起。参照文献\cite{vasconcelos2024}的联合拟合思路，我们先做一个与层数无关的选择。候选窗口取 3、5、8、12、20、30 个收益。对每个候选值，把训练段各正方差窗口视作一个局部高斯，再将这些高斯等权相加。若 $K_M^+$ 为正方差窗口数，收益直方图第 $j$ 箱的中心、密度和频数分别为 $x_j,\widehat p_j,c_j$，则
\begin{align}
 \widehat P_M(x)&=\frac{1}{K_M^+}\sum_{\widehat v_{t,M}>0}\frac{\exp[-x^2/(2\widehat v_{t,M})]}{\sqrt{2\pi\widehat v_{t,M}}},\\
 S(M)&=\frac{\sum_{j:c_j>0}c_j|\ln\widehat P_M(x_j)-\ln\widehat p_j|}{\sum_{j:c_j>0}c_j}.
\end{align}
取 $S(M)$ 最小且至少有 100 个有效窗口的 $M$。中国平安取到 $M=8$，579 个训练日、每天两段、每段 23 个收益，共有 $579\times2\times(23-8+1)=18,528$ 个窗口。以累计概率最大差作替代准则，最优窗口变为 12。这个差异说明，八个收益是本次密度估计的选择，估计所得的 40 分钟不代表一个已识别的市场物理周期。
\begin{table}[H]
\centering
\small
\caption{主例的训练期窗口选择。表中小数作四舍五入。}
\label{tab:method_windows}
\begin{tabular}{rrrr}
\toprule
$M$ & 窗口数 & $S(M)$ & 最大累计概率差\\
\midrule
3 & 24\,318 & 0.181940 & 0.054915\\
5 & 22\,002 & 0.113834 & 0.035209\\
8 & 18\,528 & 0.104956 & 0.028106\\
12 & 13\,896 & 0.107600 & 0.024120\\
20 & 4\,632 & 0.116329 & 0.024479\\
30 & 0 & 无有效候选 & 无有效候选\\
\bottomrule
\end{tabular}
\end{table}
\subsubsection{再联合拟合两种密度}
收益分箱边界取训练分位数绝对值的较大者 $a=\max(|q_{0.001}|,|q_{0.999}|)$，在 $[-a,a]$ 以初始宽度 $2\,\mathrm{IQR}\,n^{-1/3}$ 划箱，最终限制在 20--80 箱。背景在 $\ln v$ 的 0.005 至 0.995 分位间等分 30 箱。平安使用 80 个收益箱，$a=4.596$。经验密度均在各自拟合区间内归一；对数方差直方图转回 $\widehat f_k$ 时除以中心 $v_k$。给定 $N$ 和固定的 $M$，令
\begin{equation}
 E(N,\beta)=\frac{\sum_{c_j>0}c_j\left|\ln P_N(x_j)-\ln\widehat p_j\right|
 +\sum_{d_k>0}d_k\left|\ln f_N(v_k)-\ln\widehat f_k\right|}{2n-K}.\label{eq:loss}
\end{equation}
$d_k$ 是方差箱频数，$n$ 和 $K$ 是训练收益数和滚动方差窗口数。外汇论文式(24)在连续序列上按 $2n-M^*$ 归一。本文的分段实现使用 $2n-K$，主例为 34,740；同一窗口下这是常数，不影响最优形状或层数。不同组合的误差绝对值不宜直接比较。空箱跳过，对数下限为 $10^{-300}$。对每个 $N$ 分别求最小误差，才比较层数。拟合区间外的有效收益全部保留用于预测评分。

\clearpage
\subsection{深层极限、数值精度与评价期}\label{sec:oos_method}
初始训练在 1--5 层及主例的 1--8 层中搜索，约束 $\beta\in[0.05,20]$，在 $\ln\beta$ 轴作一维有界优化。若最佳点在最后一层，只能说明区间内还在改善。为检查是否会出现内部低点，对每个组合固定同一训练样本、分箱和窗口，逐层算到 32 层，随后计算 $64,128,\ldots,8192$。每层把形状上限放宽到 $\max(20,4N)$，在对数坐标取 49 个起点网格，细化局部低点并检查端点；搜索容差 $10^{-7}$。

深层模型有一个可独立计算的参照。对数方差为 $N$ 个独立 $\ln Y_i$ 之和。当 $N,\beta$ 同时增大、$N/\beta\to\theta>0$ 时，
\[
 \mathbb E[\ln v_N]=N[\ln\beta-\psi(\beta+1)]\longrightarrow-\theta/2,\qquad
 \operatorname{Var}(\ln v_N)=N\psi_1(\beta+1)\longrightarrow\theta.
\]
在 $v_0=1$ 下，极限背景为对数正态分布，$\ln v_\infty\sim\mathcal N(-\theta/2,\theta)$。把它代回式\eqref{eq:mix}，用同一联合目标在 $10^{-4}\leq\theta\leq100$ 内估计极限误差。连续两次层数翻倍的误差变化均小于 $10^{-4}$ 且相对变化小于 $0.1\%$，并且末点与极限的相对差也小于 $0.1\%$，才记为进入数值平台。若途中出现更低误差，另记下该层数。

\paragraph{积分如何落到数值上。}
在 $\ln v$ 坐标将单层乘子的密度离散为质量，用快速傅里叶卷积求和，再按式\eqref{eq:mix}混合各方差下的高斯。常规使用 $2^{16}$ 个网格点和 4,096 个混合积分点。深层时网格以理论对数均值为中心，半宽至少为 45，且覆盖 12 个对数标准差。零收益处直接用 $\mathbb E[v_N^{-1/2}]/\sqrt{2\pi}$。对最终候选再以 $2^{18}$ 个网格点、双倍区间、16,384 个积分点、97 个搜索网格点及独立的特征函数反演核对。在实际观测范围的核验点，联合误差和对数密度的路径间绝对差均小于 $2\times10^{-6}$。

\paragraph{怎样判断预测是否改善。}
每个后续收益 $x_t$ 都按观测前确定的密度评分。模型 $A$ 相对 $B$ 的平均对数评分差为
\begin{equation}
 \Delta\bar\ell=\frac{1}{n_{\rm ev}}\sum_{t=1}^{n_{\rm ev}}\ln\frac{p_{A,t}(x_t)}{p_{B,t}(x_t)}.\label{eq:score}
\end{equation}
正值表示模型 $A$ 平均给实际发生的收益更高的密度。以整交易日为配对单位，循环移动 5 日的块，重抽 2,000 次，取 2.5\% 和 97.5\% 分位数作为参考区间；另检查 1 日与 10 日块。区间条件于已估参数，各股又共同经历市场行情。2019 年测试段和 2026 年行情均在参数训练区间以外。研究设计和条件模型在历史回顾中形成，独立确认仍需要开发完成后取得的新行情。

\clearpage
\section{数据去噪与极端参数检验结果}
\subsection{先看时间顺序：前一窗口活跃，后一窗口会怎样？}
厚尾直方图在打乱交易日内的收益顺序后仍可保持原样。为了考察波动背景是否具有时间上的持续性，我们用八只股票的五分钟收益作另一项检验。每个上午、下午各取开头 16 个有效收益对应的固定时刻，按先后切成两个各含八个收益的窗口。两个窗口没有共同的收益。以 $s$ 表示上午或下午，$d$ 表示日期，分别计算两个窗口的波动强度
\[
 A_{d,s}=\tfrac18\sum_{k=1}^{8}x_{d,s,k}^2,\qquad
 B_{d,s}=\tfrac18\sum_{k=9}^{16}x_{d,s,k}^2 .
\]
八股统一取八个收益，便于在相同时间尺度上比较。上午和下午分别从 09:35、13:05 开始的五分钟区间取样。窗口使用固定钟点，不足 16 个相应位置的时段不以其他位置替代；统计仅纳入上午、下午都具备这两段窗口的日期。先在各自时段内跨日计算 $A$ 与 $B$ 的斯皮尔曼秩相关，再将上午、下午的两个系数等权平均。这样，上午开盘本身的活跃程度不会仅因和下午合并而制造相关。训练、历史测试和 2026 年分别计算；区间以日期为单位，用 5 日循环块重抽 2,000 次，每次重新计算秩，保留日内配对。该诊断描述已观测到的相关关系。
\chart{.95}{persistence_forest}{前后两个 40 分钟窗口的波动强度相关。点为上午、下午秩相关的平均值，横线为 5 日块 95\% 参考区间。两图使用同一横轴；虚线为零。华友钴业在 2026 年仅有 45 日。}{fig:persistence}
中国平安在训练段的平均相关为 0.579，2019 年测试段为 0.379（区间 0.238--0.494），2026 年为 0.450（0.323--0.548）。测试段八股点估计在 0.157--0.416 之间，其中中国建筑的区间跨零；2026 年八股点估计在 0.153--0.569 之间。相邻窗口的波动具有持续性，与慢背景影响短时收益的设想相容。单层波动过程、新闻冲击或共同市场环境也可能产生这一现象。要讨论层数，还需检验模型对两种分布的约束。

\clearpage
\subsection{从背景方差到收益密度：中国平安的训练样本}
分层模型有两个可以分别检查的环节。首先，用股票数据估计的局部方差是否与 $f_N(v)$ 接近？其次，把这个方差分布放入式\eqref{eq:mix}，能否解释五分钟收益为何既有尖峰又有厚尾？图\ref{fig:mixture}把这两步放在一起，所用的股票、训练日期及归一化尺度完全相同。
\chart{.98}{variance_mixture}{中国平安训练段的方差混合。（a）点为八个连续收益的样本方差所给出的经验密度，实线为联合拟合的单层和八层背景密度。（b）三条细线分别在训练窗口方差的第 10、50、90 百分位上画出条件高斯密度；粗线为对完整八层背景分布积分所得的收益密度。}{fig:mixture}
每个窗口只含八个收益，图\ref{fig:mixture}左侧的点本身也有抽样误差。18,528 个窗口中，有 18,342 个正方差落入作图使用的对数分箱范围。单层用一个逆伽马分布来描述这些点，最佳 $\beta$ 为 0.5627；八层用八个较平缓的随机乘子相乘，最佳 $\beta$ 为 7.9007。两组参数均同时使用收益与方差资料估计。

右图解释了混合所做的运算。三个代表性方差依次为 0.1405、0.6388 和 2.0345，单位都是训练方差。平静窗口的收益集中在零附近，剧烈窗口则给远离零点的收益留下较大的概率。模型把各种方差下的高斯密度按 $f_N(v)$ 加权。三个分位点只用来展示曲线宽度的变化，计算橙色曲线时对整个方差范围积分。于是同一组参数必须同时解释左图中的方差和实际观测的收益，后一项可在随后的测试中单独检验。

\begin{table}[H]\centering\small
\caption{中国平安静态联合误差的两部分。各行在同一训练样本、窗口与分箱上求值。}\begin{tabular}{lrrr}\toprule
层数 & 收益项 & 背景方差项 & 合计 $E$\\\midrule
单层 & 0.1282 & 0.3333 & 0.4615\\
八层 & 0.0877 & 0.0981 & 0.1858\\\bottomrule
\end{tabular}\end{table}
联合误差下降的 85.3\% 来自方差项。仅在收益密度上比较两条曲线，容易低估增加层数对背景分布的影响。反过来，训练期的方差分布拟合得更好，也并不保证下一段行情仍有相同的波动水平。


\clearpage
\subsection{静态拟合：什么变好了，什么没有被识别？}
在中国平安的训练段，八层模型同时描述了零点附近的高峰及两侧较厚的尾部。将同一组参数用于 2019 年测试段，图\ref{fig:static}右侧黑点比预测曲线更集中。该段归一化收益的标准差只有训练段的 0.5384 倍；价格变化本身已经变小。
\chart{.97}{density_plot}{中国平安五分钟收益分布。训练和测试共用分箱边界和训练尺度。每个点代表一个非空箱的密度，训练段 26,634 条、测试段 3,174 条。灰色为高斯基线，彩色为联合拟合的单层和八层模型。}{fig:static}
重抽 579 个训练日并重新选择窗口与初始层数，20 次中有 19 次选 $M=8$，20 次的层数均落在八层上限。收益中心的跨期偏差与训练误差降幅可以同时存在。图\ref{fig:depth}进一步对照了两种误差曲线。
\chart{.97}{depth_extension}{固定训练窗口后重新拟合每个层数。平安随层数增加逐渐接近极限；中信证券在九层附近有内部低点。左图横轴按 2 为底取对数，右图放大低点附近。}{fig:depth}

\clearpage
\subsection{扩大层数搜索之后}
\begin{table}[H]\centering\small
\caption{扩大层数后的训练结果。各组合固定本次训练所选窗口和分箱。“边界”表示最低值仍在已测末点，不能视为已识别的最佳层数；末点按统一稳定标准确定。15分钟为探索性对照。}\label{tab:depth_panel}
\begin{tabular}{lrrrr}\toprule
组合 & 已测最低点 & 极限核查末点 & $E_{\min}$ & $E_\infty$\\\midrule
中国平安 5分钟原始 & 4096（边界） & 4096 & 0.162176 & 0.162132\\
贵州茅台 5分钟原始 & 1024（边界） & 1024 & 0.106909 & 0.106862\\
中信证券 5分钟原始 & 9（内部） & 2048 & 0.322405 & 0.327717\\
伊利股份 5分钟原始 & 4096（边界） & 4096 & 0.168931 & 0.168897\\
招商银行 5分钟原始 & 8192（边界） & 8192 & 0.201080 & 0.201055\\
华友钴业 5分钟原始 & 7（内部） & 2048 & 0.103919 & 0.111057\\
方大炭素 5分钟原始 & 32（内部） & 2048 & 0.216678 & 0.219186\\
中国建筑 5分钟原始 & 8（内部） & 2048 & 0.641215 & 0.647902\\
中国平安 5分钟去季节性 & 4096（边界） & 4096 & 0.146240 & 0.146212\\
中国平安 5分钟剔除零收益 & 4096（边界） & 4096 & 0.147571 & 0.147543\\
中国平安 1分钟原始 & 8192（边界） & 8192 & 0.243516 & 0.243484\\
中国平安 15分钟原始 & 2048（边界） & 2048 & 0.140499 & 0.140452\\
\bottomrule\end{tabular}

\end{table}
在五分钟原始收益上，中信证券、华友钴业、方大炭素和中国建筑分别在 9、7、32 和 8 层达到已测节点的内部最低点。方大炭素还补算了 31、33 层。其余组合的最低值在末点附近，与大层数极限很接近。以平安为例，4,096 层的 $E=0.162176$，极限为 0.162132。这样的精度说明可以停止加层，报价数据仍不足以给出同样精确的市场层数。32 层之后是倍增网格，有限低点也有待在重抽训练样本后再检验。
\begin{table}[H]\centering\small
\caption{主例扩层比较的平均对数评分（纳特/条），越高越好。参数仅由训练段估计。2019年验证、测试均为短样本；检验范围见第\ref{sec:oos_method}节。}\label{tab:depth_scores}
\begin{tabular}{lrrrr}\toprule
模型 & 训练 & 2019验证 & 2019测试 & 2026\\\midrule
$N=1$ & -1.377451 & -1.052710 & -0.864814 & -1.230345\\
$N=8$ & -1.364181 & -1.074284 & -0.907113 & -1.235076\\
$N=4096$ & -1.363370 & -1.077228 & -0.912540 & -1.236412\\
大层数极限 & -1.363369 & -1.077229 & -0.912539 & -1.236412\\
\bottomrule\end{tabular}

\end{table}
在两个 2019 年留出段，单层给观测的平均密度反而较高。测试段单层减 4,096 层的评分差为 0.0477，5 日块区间为 $[0.0416,0.0538]$；2026 年差为 0.0061，区间 $[-0.0047,0.0160]$。扩大搜索解决了训练目标是否已接近极限的问题，却没有消除后续时期的尺度偏移。由此可以将改动集中到模型的最慢背景。

\clearpage
\subsection{适应股票行情：让慢背景在每日开盘前更新}
静态模型取 $v_0=1$，因为训练收益已按其标准差归一。单位可以一直沿用，市场实际波动却会改变。由 $v_N=v_0\prod_iY_i$ 可知，换一个 $v_0$ 只会按 $\sqrt{v_0}$ 拉伸收益轴。于是保留 $P_N$ 的同一形状，在开盘前根据过去的收益设置当天的尺度。
\begin{align}
 R_d&=\frac1{m_d}\sum_{k=1}^{m_d}x_{d,k}^2,\qquad
 q_{d+1}=\lambda q_d+(1-\lambda)R_d,\qquad \lambda=2^{-1/h},\label{eq:update}\\
 p_d(x\mid\mathcal F_{d-1})&=(c q_d)^{-1/2}
 P_N\!\left(x/\sqrt{c q_d};\beta,1\right).\label{eq:cond}
\end{align}
$R_d$ 为收盘后才得到的当日平均平方收益，$m_d$ 为有效收益数。$q_d$ 只用截至上一个保留交易日的收益，$h$ 为权重减半所需的观测日数。这种指数加权借鉴了风险度量中的方差更新\cite{riskmetrics1996}。$c$ 为训练所得的水平校准系数。当天每条五分钟收益使用相同的 $c q_d$；这一近似允许较快层适应当天背景，日内突发冲击仍会滞后。

\paragraph{估计和选择。}
首个交易日取 $q_1=1$，前 60 个训练日只作预热。对固定 $N,h$，最小化其后所有训练收益的负条件对数似然
\[
 -\sum_{d,k}\!\left[\ln P_N\!\left(x_{d,k}/\sqrt{c q_d};\beta,1\right)
 -\tfrac12\ln(c q_d)\right].
\]
主例的估计样本为 519 日、23,874 条收益。在对数坐标用有界拟牛顿法，$\beta\in[0.05,64]$，$c\in[0.02,50]$。八层的 $\beta$ 初值为 1、8、32，单层为 0.3、2、16；保留成功结束且似然最高的解。半衰期从 5、10、20、40、60 日中按平安八层的验证评分选择，相应评分为 $-1.0380,-1.0486,-1.0684,-1.0873,-1.0959$。之后八股共用 $h=5$，各自估计形状与尺度。最优值位于候选的短端。隔离日跳过，2019 年的保留日连续更新；2026 年在长期资料缺口后从 $q_1=1$ 重新启动。

\paragraph{四种可比较的预测。}
静态似然八层取 $q_d=1$，与条件八层使用相同的训练日、优化方法及参数个数。这对比较衡量尺度更新带来的变化。同规则高斯也逐日更新并估计 $c$，用来衡量保留厚尾形状的作用；条件单层与八层则衡量进一步加层的作用。旧静态联合八层保留作参照，其样本为完整 579 个训练日。评估按每个观测当时的条件密度取对数\cite{gneiting2007}。展示各期整体分布时，按各日有效收益数加权平均预测密度。图上的平均曲线和逐条评分各有对应的计算口径。



\paragraph{一次实际的盘前计算。}
以平安 2026 年最初两个交易日为例。1 月 5 日开盘前，从 $q_1=1$ 开始，当天的预测方差为 $c q_1=1.0471$。收盘后，46 个五分钟收益给出 $R_1=1.3702$。五日半衰期对应 $\lambda=2^{-1/5}=0.87055$，因此 1 月 6 日开盘前可算出
\[
 q_2=0.87055\times1+0.12945\times1.3702=1.04792,\qquad c q_2=1.09728.
\]
这一天实际得到 $R_2=1.92173$，高于早晨的预测。将 $R_2$ 代入同一个递推式，1 月 7 日的预测方差上升到 1.21572。即使某日低估了波动，各期评分仍使用当时开盘前给出的密度。


\clearpage
\subsection{两段行情中的逐日预测：中国平安}
在平安的训练样本上，静态似然八层给出 $\widehat\beta=11.46478,\ \widehat c=1.10064$；允许背景更新后，估计值为 $\widehat\beta=13.27908,\ \widehat c=1.04710$。两条预测曲线究竟差在哪里，可以沿交易日逐一查看。
\chart{.83}{conditional_scale}{中国平安的日波动。灰线为收盘后观测的平均平方收益 $R_d$，蓝线为开盘前给出的 $c q_d$，橙虚线为训练期的单位背景。纵轴以训练收益方差为单位。两段的纵轴范围不同。}{fig:scalepair}
历史测试段平均平方收益降到 0.2898。图\ref{fig:scalepair}左侧的蓝线也随近期波动回落，平均为 0.3287。固定背景在这段行情中把过多概率分给了较大的涨跌。2026 年的波动方向几次反转，日更新能跟随持续数日的变化，却常常错过单日尖峰。
\chart{.83}{conditional_density}{中国平安的五分钟收益。点为各箱经验密度，曲线为模型密度。条件曲线按各日收益条数加权平均。预测评分按每条收益在当日密度中的位置计算。两个面板沿用训练期的 80 箱边界。}{fig:changedensity}
\begin{table}[H]\centering\small
\caption{主例平均对数评分，单位为纳特/条，越高越好。2026年末列剔除最初20个观测日，所有模型使用相同的140日、6440条收益。}\label{tab:conditional_scores}
\begin{tabular}{lrrr}\toprule
模型 & 2019测试 & 2026全期 & 2026去前20日\\\midrule
静态联合八层 & -0.907113 & -1.235076 & -1.177017\\
静态似然八层 & -0.971858 & -1.256462 & -1.204985\\
条件八层 & -0.757054 & -1.213224 & -1.157155\\
同规则高斯 & -0.791327 & -1.280835 & -1.228232\\
\bottomrule\end{tabular}

\end{table}
在相同训练日上，加入慢背景更新后，两个时期的平均对数评分分别提高 0.2148 和 0.0432 纳特/条。剔除 2026 年启动后的前 20 日，排序依旧。

\clearpage
\subsection{八只股票：背景更新跟上了哪些变化}
\chart{.87}{eight_conditional_scale_2026}{八股在 2026 年的逐日预测。每幅图对应一只股票，灰线为 $R_d$，蓝线为 $c q_d$，橙虚线为训练期单位背景。各股按各自的训练尺度归一，纵轴范围也各不相同。华友钴业从 6 月 30 日起有 45 个观测日，其余七股各有 160 日。}{fig:eightpaths}
图\ref{fig:eightpaths}中，方大炭素在年中出现一次接近训练方差九倍的观测，蓝线随后才抬升。中国建筑的波幅小得多，数次短暂放大也较快回落。两例说明日更新所能追踪的时间尺度：波动若持续多日，前一日的收益会给次日预测提供信息；若在一个交易日内突然放大又消退，预测会落后。八股的历史测试段曲线和两个时期的密度图见八股图表补充材料。

\clearpage
\subsection{逐股评分、预测区间与零收益}
\begin{table}[H]\centering\small
\caption{条件八层相对两个对照的平均评分差，单位为纳特/条。每格首行为差值，次行为5日块95\%参考区间。}
\label{tab:conditional_eight}
\begin{tabular}{lL{3.8cm}L{3.8cm}L{3.8cm}}\toprule
股票 & 2019测试：减静态八层 & 2026：减静态八层 & 2026：减同规则高斯\\\midrule
中国平安 & $0.2148$\newline$[0.1725,0.2604]$ & $0.0432$\newline$[0.0245,0.0641]$ & $0.0676$\newline$[0.0561,0.0816]$\\
贵州茅台 & $0.0518$\newline$[0.0310,0.0731]$ & $0.0888$\newline$[0.0553,0.1265]$ & $0.1401$\newline$[0.0981,0.2062]$\\
中信证券 & $0.0493$\newline$[0.0187,0.0797]$ & $0.0915$\newline$[0.0576,0.1275]$ & $0.0985$\newline$[0.0744,0.1299]$\\
伊利股份 & $0.1444$\newline$[0.1123,0.1780]$ & $0.1850$\newline$[0.1353,0.2332]$ & $0.0744$\newline$[0.0630,0.0867]$\\
招商银行 & $0.1014$\newline$[0.0658,0.1416]$ & $0.2317$\newline$[0.1788,0.2881]$ & $0.0640$\newline$[0.0527,0.0754]$\\
华友钴业$^\dagger$ & $0.0238$\newline$[0.0016,0.0477]$ & $0.0413$\newline$[0.0126,0.0726]$ & $0.0884$\newline$[0.0421,0.1484]$\\
方大炭素 & $0.0720$\newline$[0.0407,0.1046]$ & $0.0282$\newline$[0.0104,0.0466]$ & $0.1449$\newline$[0.1013,0.1976]$\\
中国建筑 & $0.0214$\newline$[0.0078,0.0355]$ & $0.0119$\newline$[0.0047,0.0189]$ & $-0.0021$\newline$[-0.0101,0.0062]$\\
\bottomrule\end{tabular}
\vspace{3pt}\parbox{.96\linewidth}{\footnotesize 静态八层指与条件模型使用相同训练日的静态似然估计。
$\dagger$ 华友钴业2026年仅45日、2070条收益，去前20日后仅25日。其余7股全期均160日、7360条收益。
历史测试段每股为68--69日。股票受到共同市场变化影响，不能视作8次独立试验。}
\end{table}
八股在两个时期的条件八层评分均高于使用相同训练日的静态似然八层。八只股票共有交易日和市场消息，表中各行是同一规则在不同股票上的表现。进一步与同规则高斯相比，2026 年有七只股票的差值为正且区间在零以上。中国建筑的差值为 $-0.0021$，区间为 $[-0.0101,0.0062]$；这只股票的资料不足以支持增加厚尾参数。

\paragraph{预测区间能否覆盖约定比例的收益？}
按当天预测分布的 2.5\% 与 97.5\% 分位点构成区间。平安历史测试段，静态联合八层的覆盖率为 99.65\%，区间过宽；条件八层为 95.43\%，接近约定的 95\%。2026 年相应数值为 96.15\% 和 95.37\%。再以预测累计分布计算每个观测对应的分位值，其经验分布与均匀分布的最大差距在测试段从 0.1140 降至 0.0517，2026 年则从 0.0476 略升至 0.0482。改变尺度可以改善区间宽度，整条分布仍有偏差。

\paragraph{加层的贡献有多大？}
在相同日更新规则下，平安的八层减单层评分差在测试段只有 0.0009，区间 $[-0.000034,0.001726]$；2026 年为 0.0017，区间 $[0.0009,0.0024]$。八层相对单层的差值远小于从固定背景改为日更新的差值。逐笔价格有最小报价单位，相邻两个五分钟收盘价也可能完全相同。2026 年中国建筑有 46.1\% 的零收益，方大炭素为 31.6\%。任一连续密度在单点的概率都是零，因此这两只股票的分布中心即使调整了方差也仍有偏差。

\clearpage
\subsection{休市能否视为背景演化的暂停？}
五分钟收益的构造避开了午休和隔夜，但背景是否也会随闭市停下来，是另外一个可检验的问题。令开市时活动率 $a_\chi(t)=1$，闭市时 $a_\chi(t)=\chi$，其中 $0\leq\chi\leq1$。以累计活动时间 $\mathcal T_\chi(t)=\int_0^t a_\chi(s)\,\mathrm ds$ 代入随机方程，漂移和噪声的系数随之变为
\begin{align*}
 \mathrm dx&=-a_\chi\mu x\,\mathrm dt+\sqrt{2a_\chi\mu v_N}\,\mathrm dW_0,\\
 \mathrm dv_i&=-a_\chi\gamma_i(v_i-v_{i-1})\,\mathrm dt+\sqrt{a_\chi}\,\kappa_i v_i\,\mathrm dW_i.
\end{align*}
$\chi=0$ 时背景在闭市期间冻结，$\chi=1$ 时它按开市的速率演化。若各层本来处于定态，统一改变时钟只会改变两次观测之间走过的有效时间，单时点的分布保持不变。因此，只凭收益密度图无法区分这两种情形，需要考察相隔一段休市时间的两次观测。

平安历史段有 730 对相邻交易日，其中 577 对的隔夜闭市恰为 1,110 分钟。对同一组日期，分别计算午休前后和隔夜前后绝对收益的相关，两者之差记作 $D$。在交易日网格上作循环分块重抽，每次用同一批日期重算两组相关。取 5 日块的结果见表\ref{tab:market_clock_boot}。
\begin{table}[H]\centering\small
\caption{共同日期上的绝对收益皮尔逊相关。区间对应$D=\rho_{\rm overnight}-\rho_{\rm lunch}$，采用5日块。间隔1日指1110分钟闭市；全部间隔组包含周末和节假日。}\label{tab:market_clock_boot}
\begin{tabular}{lrrrr}\toprule
样本 & 配对数 & 午休相关 & 隔夜相关 & $D$的95\%区间\\\midrule
全期，全部间隔 & 730 & 0.1850 & 0.1542 & $[-0.1402,0.0790]$\\
全期，间隔1日 & 577 & 0.2077 & 0.1799 & $[-0.1545,0.1075]$\\
测试期，全部间隔 & 68 & 0.0148 & 0.0253 & $[-0.4168,0.4167]$\\
测试期，间隔1日 & 53 & 0.0232 & 0.0688 & $[-0.3755,0.4236]$\\
\bottomrule\end{tabular}

\end{table}
全期及测试段的区间均跨过零。仅看固定长度的隔夜间隔，测试段只剩 53 对观测。以这批数据还难以区分休市时间、盘中节律和隔夜消息各自的作用，故未将 $\chi$ 加入逐日预测的待估参数。

\subsection{结论}
湍流类比在这里给出了一个可操作的问题：如果收益受缓慢变化的波动背景调节，那么，相邻但互不重叠的窗口是否会一起活跃？八股的样本给出了多数为正且可辨别的关系。对随后一个五分钟区间，层级混合提供了收益可能落在哪里的分布；单凭边际密度却难以分清许多相邻层数。把当天的背景尺度交给过去交易日的波动来估计，对八股的后续评分产生了比平安单纯加层更清楚的改善。

这一结果适用于本文选取的活跃沪市股票及两段观测期。历史测试段每股仅 68 至 69 日，2026 年华友钴业仅 45 日，表中的区间也以已估参数为条件。若要判断日内尖峰或零收益的发生概率，还需要能在日内更新的状态、对报价离散性的建模，以及新的跨期检验。

\clearpage
\noindent\begin{minipage}[t][0.5\textheight][t]{\textwidth}
\section{AI 辅助诊断专栏}\label{sec:ai_audit}
AI用于辅助编程、推导核对与文字整理。以下保留历史开发记录中的两项数值问题，
说明如何通过概率守恒、坐标变换和独立积分检查计算，而非仅凭曲线外观判断正确性。

\paragraph{归一化因子的漏乘。}
开发记录中的旧密度实现把各层的伽马函数因子乘积写成了单个因子。正确分母为
$\prod_{i=1}^N\Gamma(\beta_i+1)$，共同$\beta$时为$\Gamma(\beta+1)^N$。
漏乘不改变曲线形状，却使背景密度整体多出$\Gamma(\beta+1)^{N-1}$倍。
以$N=2,\beta=1.5$逐层独立积分，错误实现与积分结果之比恒为
$\Gamma(2.5)=1.3293$，据此定位漏项。修正后，112个收益测试点中有1点参考求值失败，
其余点的最大相对差为$2.0\times10^{-14}$；失败点不记作通过。


\paragraph{卷积坐标与雅可比错误。}
旧卷积实现把数组首项当作对数坐标原点，造成结果平移；从$y=\ln v$的密度$g(y)$
还原为方差密度时，又漏掉了$1/v$。概率守恒要求
$f(v)\,\mathrm dv=g(y)\,\mathrm dy$，故$f(v)=g(\ln v)/v$。
与直接积分比较时，漏掉雅可比的结果恰好多出$v$倍。
修正原点移位和雅可比后，84个背景测试点均完成求值；其中2点密度低于$10^{-10}$，
不作相对误差比较，其余点最大相对差为$7.5\times10^{-6}$。
扩层计算另用加密网格与特征函数反演检查实际观测范围内的精度。

\paragraph{公式与文字的一致性。}
AI草稿曾将单层$t$基线称为“固定尺度”。实际上固定的是$v_0=1$，
尺度平方$\beta v_0/(\beta+1)$随$\beta$变化。现已改为“固定位置和背景均值，仅估计$\beta$”，
并明确其与三参数$t$的区别；这处修正不改变已经计算的评分。
上述复核包含AI辅助步骤，不能替代成员对公式和核心函数的独立解释。
\end{minipage}

\section{成员贡献度量化声明}\label{sec:contributions}
\begin{center}
\setlength{\tabcolsep}{20pt}
\begin{tabular}{lcc}\toprule
成员 & 算法 & 汇报\\\midrule
方洪威 & 70\% & 30\%\\
戴晴琴 & 30\% & 70\%\\\bottomrule
\end{tabular}
\end{center}



\paragraph{数据与材料可用性。}
两个研究时期的八股原始一分钟行情共 7,009 份逐日文件、1,678,159 条记录。仓库同时提供五分钟收益、完整股票池的筛选统计、拟合程序、逐条评分和不重叠窗口实验的输入说明与结果。完整选股表和八股曲线分别见两份补充材料。
\small\begin{thebibliography}{9}\setlength{\itemsep}{0pt}
\bibitem{vasconcelos2024} G. L. Vasconcelos, L. R. C. Ribeiro, A. M. S. Mac\^edo,
I. R. R. Gonz\'alez, R. Ospina, and A. A. Brum, Turbulence hierarchy in foreign exchange markets,
\emph{Phys. Rev. E} \textbf{109}, 044313 (2024).
\bibitem{macedo2017} A. M. S. Mac\^edo, I. R. R. Gonz\'alez, D. S. P. Salazar, and G. L. Vasconcelos,
Universality classes of fluctuation dynamics in hierarchical complex systems,
\emph{Phys. Rev. E} \textbf{95}, 032315 (2017).
\bibitem{riskmetrics1996} J. P. Morgan and Reuters, \emph{RiskMetrics---Technical Document},
4th ed. (New York, 1996), Sec. 5.2.
\bibitem{gneiting2007} T. Gneiting and A. E. Raftery, Strictly proper scoring rules, prediction,
and estimation, \emph{J. Am. Stat. Assoc.} \textbf{102}, 359--378 (2007).
\end{thebibliography}
\par\smallskip
\noindent\textbf{项目仓库}\quad
{\small\url{https://github.com/fanghongwei1017/PJ5-Hierarchical-Volatility}}
\end{document}
